\section*{الفصل \ref{ch:g3:mult} --- الضرب}

\begin{solution}{exo:g3:mult:1}
$8 + 8 + 8 + 8 = 4 \times 8 = 32$؛ \quad
$5 + 5 + 5 = 3 \times 5 = 15$؛ \quad
$20 \times 5 = 100$.
\end{solution}

\begin{solution}{exo:g3:mult:2}
$42$؛ \quad $32$؛ \quad $81$؛ \quad $56$؛ \quad $36$.
\end{solution}

\begin{solution}{exo:g3:mult:3}
المستطيلان يحتويان النقط $24$ نفسها --- وأحدهما هو الآخر
مُدارًا على جنبه: $4 \times 6 = 6 \times 4$.
\end{solution}

\begin{solution}{exo:g3:mult:4}
$560$؛ \quad $800$؛ \quad $0$؛ \quad $45$؛ \quad $300$.
\end{solution}

\begin{solution}{exo:g3:mult:5}
$132 \times 3 = 396$؛ \quad
$217 \times 4 = 868$؛ \quad
$408 \times 7 = 2\,856$.
\end{solution}

\begin{solution}{exo:g3:mult:6}
$31 \times 5 = 30 \times 5 + 5 = 155$؛ \quad
$42 \times 3 = 120 + 6 = 126$؛ \quad
$25 \times 6 = 20 \times 6 + 5 \times 6 = 120 + 30 = 150$.
\end{solution}

\begin{solution}{exo:g3:mult:7}
$99 \times 7 = 700 - 7 = 693$؛ \quad
$101 \times 6 = 600 + 6 = 606$؛ \quad
$98 \times 5 = 500 - 10 = 490$.
\end{solution}

\begin{solution}{exo:g3:mult:8}
$389 \times 5$: التقدير $400 \times 5 = 2\,000$؛ وبالضبط
$1\,945$.

$612 \times 8$: التقدير $600 \times 8 = 4\,800$؛ وبالضبط $4\,896$.
\end{solution}

\begin{solution}{exo:g3:mult:9}
$38 \times 6 = 228$. في العلب $228$ بيضة.
\end{solution}

\begin{solution}{exo:g3:mult:10}
الطاولات: $8 \times 4 = 32$. والتلاميذ: $32 \times 2 = 64$. فتتسع
الغرفة $64$ تلميذًا.
\end{solution}

\begin{solution}{exo:g3:mult:11}
يظهر $24$ في مربّع الجداول على صور $1 \times 24$ و $2 \times 12$
و $3 \times 8$ و $4 \times 6$ --- والأزواج نفسها معكوسة. والترتيبات
الممكنة: $1$ صف من $24$، و $2$ من الصفوف من $12$، و $3$ صفوف من $8$،
و $4$ صفوف من $6$، و $6$ صفوف من $4$، و $8$ صفوف من $3$، و $12$ صفًّا من $2$،
و $24$ صفًّا من $1$.
\end{solution}

\begin{solution}{exo:g3:mult:12}
تفكّك زينة $47$ إلى $40 + 7$ (وهي طريقة الأعمدة ذهنيًّا)؛ أما
لينا فتستبدل بالعدد $47$ العددَ الودود $50$ ثم تصحّح بطرح
$3 \times 6$ التي عُدّت زيادةً. وفي $38 \times 4$: بطريقة زينة،
$120 + 32 = 152$؛ وبطريقة لينا، $40 \times 4 - 2 \times 4 = 160 - 8 =
152$. وكلتاهما تنجح --- وطريقة لينا أسرع هنا لأن $38$ قريب من
$40$.
\end{solution}
